\begin{titlepage}
	\begin{tikzpicture}[remember picture, overlay]
		% ---- 1. 灰色渐变铺满整个封面（金属灰 → 深灰黑）----
		\shade[top color=DarkBlue, middle color=DarkBlue!90, bottom color=DarkBlue!65!black]
			(current page.south west) rectangle (current page.north east);

		% ---- 2. 细点阵（模拟网格，微弱纹理）----
		\foreach \x in {1,...,20} {
			\foreach \y in {1,...,29} {
				\fill[white, opacity=0.05] (\x,\y) circle[radius=0.018cm];
			}
		}

		% ---- 3. 主标题后方的柔光晕 ----
		\fill[white, opacity=0.05] ([yshift=-8.2cm]current page.north) circle[radius=7.0cm];
		\fill[white, opacity=0.05] ([yshift=-8.2cm]current page.north) circle[radius=4.6cm];

		% ---- 4. 右上：材料拉伸应力—应变曲线（σ-ε 曲线）----
		\begin{scope}[shift={([xshift=-6.2cm,yshift=-6.2cm]current page.north east)}]
			% 坐标轴
			\draw[white, opacity=0.55, line width=0.8pt, -{Stealth}] (0,0) -- (0,5.4);
			\draw[white, opacity=0.55, line width=0.8pt, -{Stealth}] (0,0) -- (6.2,0);
			% 曲线：线弹性段 + 屈服平台 + 强化段
			\draw[white, line width=1.1pt, opacity=0.9, domain=0:2.0, samples=60, smooth]
				plot (\x, {2.1*\x});
			\draw[white, line width=1.1pt, opacity=0.9] (2.0,4.2) -- (2.7,4.2);
			\draw[white, line width=1.1pt, opacity=0.9, domain=2.7:5.4, samples=80, smooth]
				plot (\x, {4.2 + 1.0*(1-exp(-0.8*(\x-2.7)))});
			% 屈服点辅助虚线
			\draw[white, dashed, opacity=0.4, line width=0.7pt] (2.0,0) -- (2.0,4.2);
			\draw[white, dashed, opacity=0.4, line width=0.7pt] (0,4.2) -- (2.0,4.2);
			% 横轴刻度
			\foreach \i in {1.0,2.0,3.0,4.0,5.0,6.0} {
				\draw[white, opacity=0.3, line width=0.5pt] (\i,-0.08) -- (\i,0.08);
			}
			\node[anchor=north, text=white, opacity=0.7, font=\small\itshape] at (6.2,-0.3) {$\varepsilon$};
			\node[anchor=south east, text=white, opacity=0.7, font=\small\itshape] at (-0.15,5.4) {$\sigma$};
			\node[anchor=south east, text=white, opacity=0.8, font=\small\itshape] at (6.0,-3.6) {应力—应变曲线};
		\end{scope}

		% ---- 5. 左下：简支梁弯曲变形（集中力 + 挠曲线）----
		\begin{scope}[shift={([xshift=1.0cm,yshift=1.1cm]current page.south west)}, scale=1.25]
			% 梁
			\draw[white, opacity=0.95, line width=1.6pt] (0.5,0.9) -- (5.5,0.9);
			% 左端固定铰支座（V 形 + 地面剖面线）
			\draw[white, line width=1.2pt, opacity=0.9] (0.5,0.9) -- (0.78,0.62) -- (1.06,0.9);
			\draw[white, line width=1.4pt, opacity=0.85] (0.2,0.62) -- (1.36,0.62);
			\foreach \x in {0.2,0.45,...,1.2} {
				\draw[white, opacity=0.4, line width=0.7pt] (\x,0.62) -- (\x-0.15,0.34);
			}
			% 右端滚动支座（三角形滚轮）
			\draw[white, line width=1.2pt, opacity=0.9] (5.5,0.9) -- (5.78,0.62) -- (5.22,0.62) -- cycle;
			\draw[white, line width=1.4pt, opacity=0.85] (4.9,0.62) -- (6.1,0.62);
			% 集中力
			\draw[white, line width=1.4pt, opacity=0.9, -{Stealth}] (3.0,3.3) -- (3.0,1.25);
			\node[anchor=south, text=white, opacity=0.8, font=\small\itshape] at (3.0,3.4) {$F$};
			% 挠曲线（虚线）
			\draw[white, dashed, opacity=0.7, line width=1.0pt, domain=0.5:5.5, samples=80, smooth]
				plot (\x, {0.9 - 0.55*sin(deg(pi*(\x-0.5)/5))});
			% 标签
			\node[anchor=north west, text=white, opacity=0.8, font=\small\itshape] at (0,3.9) {简支梁弯曲变形};
		\end{scope}

		% ---- 6. 银色细边框 ----
		\draw[gold, line width=0.7pt, opacity=0.55]
			([shift={(0.85,0.85)}]current page.south west) rectangle
			([shift={(-0.85,-0.85)}]current page.north east);

		% ---- 7. 顶部眉题（英文小字 + 银色短线）----
		\coordinate (T) at ([yshift=-3.5cm]current page.north);
		\draw[gold, line width=1.0pt] (T) ++(-3.4cm,0) -- ++(1.0cm,0);
		\draw[gold, line width=1.0pt] (T) ++(2.4cm,0) -- ++(1.0cm,0);
		\node[anchor=north, text=white, opacity=0.85, align=center] at (T) {
			{\sffamily\fontsize{13}{16}\selectfont
				\uppercase{Course \ Notes}}
		};

		% ---- 8. 主标题（居中，使用 \notename）----
		\node[anchor=center, align=center, text=white] at ([yshift=-8.0cm]current page.north) {%
			{\fontsize{40}{50}\sffamily\bfseries \notename}
		};

		% ---- 9. 银色菱形 + 双细线 ----
		\coordinate (C) at ([yshift=-11.4cm]current page.north);
		\draw[gold, line width=1.1pt] (C) ++(-2.2cm,0) -- ++(1.55cm,0);
		\draw[gold, line width=1.1pt] (C) ++(0.65cm,0) -- ++(1.55cm,0);
		\node[fill=gold, minimum size=0.24cm, inner sep=0pt, rotate=45] at (C) {};

		% ---- 10. 副标题 ----
		\node[anchor=north, text=white, opacity=0.9, align=center]
			at ([yshift=-12.2cm]current page.north) {
			{\fontsize{16}{20}\sffamily 学\ \ 习\ \ 笔\ \ 记}
		};

		% ---- 11. 底部作者、日期信息 ----
		\coordinate (B) at ([yshift=3.4cm]current page.south);
		\draw[gold, line width=0.8pt, opacity=0.7] (B) ++(-1.1cm,0.85cm) -- ++(2.2cm,0);
		\node[anchor=south, text=white, align=center] at (B) {
			\begin{tabular}{c}
				{\fontsize{20}{24}\sffamily\bfseries \noteauthor} \\[0.4cm]
				{\large \sffamily \today}
			\end{tabular}
		};
	\end{tikzpicture}
\end{titlepage}
